\frfilename{mt281.tex} 
\versiondate{4.12.12} 
\copyrightdate{1994} 
      
\def\chaptername{Analisis Fourier} 
\def\sectionname{Teorema Stone-Weierstrass} 
      
\newsection{281} 
      
Sebelum mulai mengerjakan pokok bahasan sesungguhnya dari bab ini, akan
bermanfaat untuk memiliki pernyataan yang cukup umum tentang suatu teorema
mendasar mengenai hampiran fungsi kontinu. Sebenarnya saya memberikan
berbagai bentuk (281A, 281E, 281F, dan 281G, bersama dengan 281Ya, 281Yd,
dan 281Yg), yang semuanya kadang-kadang berguna. Saya mengakhiri seksi ini
dengan suatu versi teorema ekuidistribusi Weyl (281M--281N).
      
\leader{281A}{Teorema Stone-Weierstrass: bentuk pertama} Misalkan $X$ adalah
ruang topologis dan $K$ adalah subhimpunan kompak dari $X$. Tuliskan
$C_b(X)$ untuk ruang semua fungsi kontinu bernilai real yang terbatas pada
$X$, sehingga $C_b(X)$ adalah ruang linear atas $\Bbb R$.
Misalkan $A\subseteq C_b(X)$ sedemikian sehingga
      
\inset{$A$ adalah subruang linear dari $C_b(X)$;}
      
\inset{$|f|\in A$ untuk setiap $f\in A$;}
      
\inset{$\chi X\in A$;}
      
\inset{untuk setiap titik berbeda $x$, $y$ dari $K$ terdapat suatu $f\in A$
sedemikian sehingga $f(x)\ne f(y)$.}
      
\noindent Maka untuk setiap $h:K\to\Bbb R$ kontinu dan $\epsilon>0$
terdapat $f\in A$ sedemikian sehingga
      
\inset{$|f(x)-h(x)|\le\epsilon$ untuk setiap $x\in K$,}
      
\inset{jika $K\ne\emptyset$, $\inf_{x\in X}f(x)\ge\inf_{x\in K}h(x)$ dan
$\sup_{x\in X}f(x)\le\sup_{x\in K}h(x)$.}
      
\cmmnt{\medskip
      
\noindent{\bf Catatan} Saya menyatakan teorema ini dalam konteks alaminya,
yaitu ruang topologis umum. Namun jika ruang-ruang ini masih asing bagi
Anda, sebenarnya Anda tidak perlu mengetahui apa itu. Jika Anda membaca
`misalkan $X$ ruang topologis' sebagai `misalkan $X$ suatu subhimpunan dari
$\Bbb R^r$' dan `$K$ adalah subhimpunan kompak dari $X$' sebagai `$K$ adalah
subhimpunan $X$ yang tertutup dan terbatas dalam $\BbbR^r$', hal itu sudah
cukup untuk semua penerapan dalam bab ini. Tentu saja, untuk mengikuti
buktinya Anda perlu mengetahui sedikit tentang kekompakan dalam $\BbbR^r$;
saya telah menuliskan fakta-fakta yang diperlukan dalam \S2A2.
}
      
\proof{{\bf (a)} Jika $K$ kosong, kita dapat mengambil $f=\tbf{0}$ sebagai
fungsi konstan bernilai $0$. Jadi selanjutnya anggaplah $K\ne\emptyset$.
      
\medskip
      
{\bf (b)} Hal pertama yang perlu diperhatikan adalah bahwa jika $f$, $g\in A$
maka $f\wedge g$ dan $f\vee g$ termasuk dalam $A$, dengan
      
\Centerline{$(f\wedge g)(x)=\min(f(x),g(x))$,
\quad $(f\vee g)(x)=\max(f(x),g(x))$}
      
\noindent untuk setiap $x\in X$; ini karena
      
\Centerline{$f\wedge g=\Bover12(f+g-|f-g|)$,
\quad $f\vee g=\Bover12(f+g+|f-g|)$.}
      
\noindent Dengan induksi pada $n$, diperoleh bahwa $f_0\wedge\ldots\wedge f_n$
dan $f_0\vee\ldots\vee f_n$ termasuk dalam $A$ untuk semua
$f_0,\ldots,f_n\in A$.
      
\medskip
      
{\bf (c)} Jika $x$, $y$ adalah titik berbeda dari $K$, dan $a$, $b\in\Bbb R$,
terdapat $f\in A$ sedemikian sehingga $f(x)=a$ dan $f(y)=b$.
\Prf\ Mulailah dengan $g\in A$ sedemikian sehingga $g(x)\ne g(y)$; di sinilah
kita menggunakan hipotesis terakhir dari empat hipotesis pada $A$. Tetapkan
      
\Centerline{$\alpha=\Bover{a-b}{g(x)-g(y)}$,\quad
$\beta=\Bover{bg(x)-ag(y)}{g(x)-g(y)}$,\quad
$f=\alpha g+\beta\chi X\in A$.   \Qed}
      
\medskip
      
{\bf (d)} (Inti bukti terletak pada dua paragraf berikutnya.) Misalkan
$h:K\to\coint{0,\infty}$ adalah fungsi kontinu dan $x$ sebarang titik dari
$K$. Untuk setiap $\epsilon>0$, terdapat $f\in A$ sedemikian sehingga
$f(x)=h(x)$ dan $f(y)\le h(y)+\epsilon$ untuk setiap $y\in K$.
\Prf\ Misalkan $\Cal G_x$ keluarga semua himpunan terbuka $G\subseteq X$
yang untuknya terdapat $f\in A$ sedemikian sehingga $f(x)=h(x)$ dan
$f(w)\le h(w)+\epsilon$ untuk setiap $w\in K\cap G$. Saya mengklaim bahwa
$K\subseteq\bigcup\Cal G_x$. Untuk melihatnya, ambil sebarang $y\in K$.
Menurut (c), terdapat $f\in A$ sedemikian sehingga $f(x)=h(x)$ dan
$f(y)=h(y)$. Kini $h-f\restr K:K\to\Bbb R$ adalah fungsi kontinu yang
bernilai $0$ di $y$, sehingga terdapat subhimpunan terbuka $G$ dari $X$ yang
memuat $y$ dan memenuhi $(h-f\restr K)(w)\ge-\epsilon$ untuk setiap
$w\in G\cap K$, yaitu $f(w)\le h(w)+\epsilon$ untuk setiap $w\in G\cap K$.
Maka $G\in\Cal G_x$ dan $y\in\bigcup\Cal G_x$, sebagaimana diperlukan.
      
Karena $K$ kompak, $\Cal G_x$ memiliki subselimut berhingga, misalnya
$G_0,\ldots,G_n$. Untuk setiap $i\le n$, ambil $f_i\in A$ sedemikian
sehingga $f_i(x)=h(x)$ dan $f_i(w)\le h(w)+\epsilon$ untuk setiap
$w\in G_i\cap K$. Maka
      
\Centerline{$f=f_0\wedge f_1\wedge\ldots\wedge f_n\in A$,}
      
\noindent menurut (b), dan jelas $f(x)=h(x)$, sedangkan jika $y\in K$ ada
$i\le n$ sedemikian sehingga $y\in G_i$, sehingga
      
\Centerline{$f(y)\le f_i(y)\le h(y)+\epsilon$.   \Qed}
      
      
\medskip
      
{\bf (e)} Jika $h:K\to\Bbb R$ adalah fungsi kontinu sebarang dan $\epsilon>0$,
terdapat $f\in A$ sedemikian sehingga $|f(y)-h(y)|\le\epsilon$ untuk setiap
$y\in K$. \Prf\ Kali ini, misalkan $\Cal G$ adalah himpunan semua
subhimpunan terbuka $G$ dari $X$ yang untuknya terdapat $f\in A$ sedemikian
sehingga $f(y)\le h(y)+\epsilon$ untuk setiap $y\in K$ dan
$f(x)\ge h(x)-\epsilon$ untuk setiap $x\in G\cap K$. Sekali lagi, $\Cal G$
adalah selimut terbuka bagi $K$. Untuk melihatnya, ambil sebarang $x\in K$.
Menurut (d), terdapat $f\in A$ sedemikian sehingga $f(x)=h(x)$ dan
$f(y)\le h(y)+\epsilon$ untuk setiap $y\in K$. Kini
$h-f\restr K:K\to\Bbb R$ adalah fungsi kontinu yang bernilai nol di $x$,
sehingga terdapat subhimpunan terbuka $G$ dari $X$, memuat $x$, sedemikian
sehingga $(h-f\restr K)(w)\le\epsilon$ untuk setiap $w\in G\cap K$, yaitu
$f(w)\ge h(w)-\epsilon$ untuk setiap $w\in G\cap K$. Maka $G\in\Cal G$
dan $x\in\bigcup\Cal G$, sebagaimana diperlukan.
      
Karena $K$ kompak, $\Cal G$ memiliki subselimut berhingga, misalnya
$G_0,\ldots,G_m$. Untuk setiap $j\le m$, ambil $f_j\in A$ sedemikian
sehingga $f_j(y)\le h(y)+\epsilon$ untuk setiap $y\in K$ dan
$f_j(w)\ge h(w)-\epsilon$ untuk setiap $w\in G_j\cap K$. Maka
      
\Centerline{$f=f_0\vee f_1\vee\ldots\vee f_m\in A$,}
      
\noindent menurut (b), dan jelas $f(y)\le h(y)+\epsilon$ untuk setiap
$y\in K$, sedangkan jika $x\in K$ ada $j\le m$ sedemikian sehingga
$x\in G_j$, sehingga
      
\Centerline{$f(x)\ge f_j(x)\ge h(x)-\epsilon$.}
      
\noindent Jadi $|f(x)-h(x)|\le\epsilon$ untuk setiap $x\in K$, sebagaimana
diperlukan. \Qed
      
\medskip
      
{\bf (f)} Jadi kita memiliki $f$ yang memenuhi persyaratan pertama dari dua
persyaratan teorema. Untuk persyaratan kedua, tetapkan
$M_0=\inf_{x\in K}h(x)$ dan $M_1=\sup_{x\in K}h(x)$, serta
      
\Centerline{$f_1=\med(M_0\chi X,f,M_1\chi X)
=(M_0\chi X)\vee(f\wedge M_1\chi X)$;}
      
\noindent $f_1$ juga memenuhi kondisi kedua selain kondisi pertama.
(Di sini saya diam-diam mengasumsikan sesuatu yang memang benar, bahwa $M_0$
dan $M_1$ berhingga; ini karena $K$ kompak -- lihat 2A2G atau 2A3N.)
}%end of proof of 281A
\leader{281B}{}\cmmnt{ Kita memerlukan beberapa alat sederhana yang termasuk
dalam teori dasar ruang bernorma; tetapi saya berharap alat-alat ini tetap
dapat diikuti bahkan jika Anda belum pernah menjumpai `ruang bernorma',
asalkan Anda menyimpan jari pada awal \S2A4 ketika membaca lemma berikutnya.
      
\medskip
      
\noindent}{\bf Lemma} Misalkan $X$ sebarang himpunan. Tuliskan
$\ell^{\infty}(X)$ untuk himpunan fungsi terbatas dari $X$ ke $\Bbb R$. Untuk
$f\in\ell^{\infty}(X)$, tetapkan
      
\Centerline{$\|f\|_{\infty}=\sup_{x\in X}|f(x)|$,}
      
\noindent dengan supremum dihitung sebagai $0$ jika $X$ kosong. Maka
      
(a) $\ell^{\infty}(X)$ adalah ruang bernorma.
      
(b) Misalkan $A\subseteq\ell^{\infty}(X)$ suatu himpunan dan
$\overline{A}$ penutupannya\cmmnt{ (2A3D)}.
      
\quad (i) Jika $A$ adalah subruang linear dari $\ell^{\infty}(X)$, demikian
juga $\overline{A}$.
      
\quad (ii) Jika $f\times g\in A$ setiap kali $f$, $g\in A$, maka
$f\times g\in\overline{A}$ setiap kali $f$, $g\in\overline{A}$.
      
\quad (iii) Jika $|f|\in A$ setiap kali $f\in A$, maka $|f|\in\overline{A}$
setiap kali $f\in\overline{A}$.
      
\proof{{\bf (a)} Ini adalah verifikasi rutin. Untuk memastikan bahwa
$\ell^{\infty}(X)$ merupakan ruang linear atas $\Bbb R$, kita harus
memeriksa bahwa $f+g$, $cf$ termasuk dalam $\ell^{\infty}(X)$ setiap kali
$f$, $g\in\ell^{\infty}(X)$ dan $c\in\Bbb R$; sekaligus kita dapat
memastikan bahwa $\|\,\|_{\infty}$ merupakan norma pada
$\ell^{\infty}(X)$ dengan mengamati bahwa
      
\Centerline{$|(f+g)(x)|\le|f(x)|+|g(x)|
\le\|f\|_{\infty}+\|g\|_{\infty}$,}
      
\Centerline{$|cf(x)|=|c||f(x)|\le|c|\|f\|_{\infty}$}
      
\noindent setiap kali $f$, $g\in\ell^{\infty}(X)$ dan $c\in\Bbb R$.
Pada saat yang sama patut dicatat bahwa jika $f$, $g\in\ell^{\infty}(X)$,
maka
      
\Centerline{$|(f\times g)(x)|
=|f(x)||g(x)|\le\|f\|_{\infty}\|g\|_{\infty}$}
      
\noindent untuk setiap $x\in X$, sehingga
$\|f\times g\|_{\infty}\le\|f\|_{\infty}\|g\|_{\infty}$.
      
(Tentu saja semua pengamatan ini merupakan kasus khusus yang sangat
elementer dari beberapa bagian \S243; lihat 243Xl.)
      
\medskip
      
{\bf (b)} Ingat bahwa
      
\Centerline{$\overline{A}=\{f:f\in\ell^{\infty}(X),\Forall
\epsilon>0\Exists f_1\in A,
\,\|f-f_1\|_{\infty}\le\epsilon\}$}
      
\noindent (2A3Kb). Ambil $f$, $g\in\overline{A}$ dan $c\in\Bbb R$, serta
misalkan $\epsilon>0$. Tetapkan
      
\Centerline{$\eta
=\min(1,\Bover{\epsilon}{2+|c|+\|f\|_{\infty}+\|g\|_{\infty}})>0$.}
      
\noindent Maka terdapat $f_1$, $g_1\in A$ sedemikian sehingga
$\|f-f_1\|_{\infty}\le\eta$ dan $\|g-g_1\|_{\infty}\le\eta$.
      
Sekarang
      
$$\eqalignno{\|(f+g)-(f_1+g_1)\|_{\infty}
&\le\|f-f_1\|_{\infty}+\|g-g_1\|_{\infty}
\le 2\eta
\le\epsilon,\cr
\cr
\|cf-cf_1\|_{\infty}
&=|c|\|f-f_1\|_{\infty}
\le|c|\eta
\le\epsilon,\cr
\cr
\|(f\times g)-(f_1\times g_1)\|_{\infty}
&=\|(f-f_1)\times g+f\times(g-g_1)-(f-f_1)\times(g-g_1)\|_{\infty}\cr
&\le\|(f-f_1)\times g\|_{\infty}+\|f\times(g-g_1)\|_{\infty}
  +\|(f-f_1)\times(g-g_1)\|_{\infty}\cr
&\le\|f-f_1\|_{\infty}\|g\|_{\infty}+\|f\|_{\infty}\|g-g_1)\|_{\infty}
  +\|f-f_1\|_{\infty}\|g-g_1\|_{\infty}\cr
&\le\eta(\|g\|_{\infty}+\|f\|_{\infty}+\eta)
\le\eta(\|g\|_{\infty}+\|f\|_{\infty}+1)
\le\epsilon,\cr
\cr
\||f|-|f_1|\|_{\infty}
&\le\|f-f_1\|_{\infty}
\le\eta
\le\epsilon.\cr}$$
      
\quad{\bf (i)} Jika $A$ adalah subruang linear, maka $f_1+g_1$ dan $cf_1$
termasuk dalam $A$. Karena $\epsilon$ sebarang, $f+g$ dan $cf$ termasuk
dalam $\overline{A}$. Karena $f$, $g$, dan $c$ sebarang, $\overline{A}$ adalah
subruang linear dari $\ell^{\infty}(X)$.
      
\quad{\bf (ii)} Jika $A$ tertutup terhadap perkalian, maka $f_1\times g_1\in A$.
Karena $\epsilon$ sebarang, $f\times g\in\overline{A}$.
      
\quad{\bf (iii)} Jika nilai mutlak fungsi-fungsi dalam $A$ termasuk dalam $A$,
maka $|f_1|\in A$. Karena $\epsilon$ sebarang, $|f|\in\overline{A}$.
}%end of proof of 281B
\leader{281C}{Lemma} Terdapat suatu barisan $\sequencen{p_n}$ polinom real
sedemikian sehingga $\lim_{n\to\infty}p_n(x)
\ifdim\pagewidth=390pt\break\fi
=|x|$ secara seragam untuk $x\in[-1,1]$.
      
\proof{{\bf (a)} Menurut Teorema Binomial, kita memiliki
      
\Centerline{$(1-x)^{1/2}
=1-\Bover12x-\Bover1{4{\cdot}2!}x^2
  -\Bover{1{\cdot}3}{2^3{\cdot}3!}x^3-\ldots
=-\sum_{n=0}^{\infty}\Bover{(2n)!}{(2n-1)(2^nn!)^2}x^n$}
      
\noindent setiap kali $|x|<1$, dengan konvergensi seragam pada setiap
interval $[-a,a]$ dengan $0\le a<1$. (Untuk bukti, lihat hampir semua buku
analisis real atau kompleks. Jika Anda tidak memiliki teks favorit, Anda
dapat mencoba menyusun bukti dari fakta-fakta berikut: (i) radius konvergensi
deret adalah $1$, sehingga pada setiap interval $[-a,a]$, dengan $0\le a<1$,
deret tersebut dapat dijumlahkan secara mutlak dan seragam (ii) dengan
menuliskan $f(x)$ untuk jumlah deret ketika $|x|<1$, gunakan Teorema
Konvergensi Terdominasi Lebesgue untuk mencari bentuk integral tak tentu
$\int_0^xf$, $-\int_{-x}^0f$ dan tunjukkan bahwa keduanya adalah
$\bover23(1-(1-x)f(x))$, $\bover23(1-(1+x)f(-x))$ untuk $0\le x<1$
(iii) gunakan Teorema Dasar Kalkulus untuk menunjukkan bahwa
$f(x)+2(1-x)f'(x)=0$
(iv) tunjukkan bahwa $\bover{d}{dx}\bigl(\bover{f(x)^2}{1-x}\bigr)=0$ dan
karena itu (v) $f(x)^2=1-x$ setiap kali $|x|<1$. Terakhir, tunjukkan bahwa
karena $f$ kontinu dan tidak nol di $\ooint{-1,1}$, $f(x)$ harus merupakan
akar kuadrat positif dari $1-x$ di seluruh interval.)
      
Kita memiliki potongan informasi tambahan. Jika kita menetapkan
      
\Centerline{$q_0(x)=1$, \quad$q_1(x)=1-\Bover12x$,
\quad$q_n(x)=-\sum_{k=0}^n\Bover{(2k)!}{(2k-1)(2^kk!)^2}x^k$}
      
\noindent untuk $n\ge 2$ dan $x\in[0,1]$, sehingga $q_n$ adalah jumlah parsial
ke-$n$ dari deret binomial untuk $(1-x)^{1/2}$, maka
$\lim_{n\to\infty}q_n(x)=(1-x)^{1/2}$ untuk setiap $x\in\coint{0,1}$.
Namun setiap $q_n$ juga tak menaik pada $[0,1]$, dan $\sequencen{q_n(x)}$
adalah barisan tak menaik untuk setiap $x\in[0,1]$. Jadi kita harus memiliki
      
\Centerline{$\sqrt{1-x}\le q_n(x)
\Forall n\in\Bbb N,\,x\in\coint{0,1}$,}
      
\noindent dan karena itu, sebab semua $q_n$ kontinu,
      
\Centerline{$\sqrt{1-x}\le q_n(x)
\Forall n\in\Bbb N,\,x\in[0,1]$.}
      
\noindent Lebih lanjut, diberikan $\epsilon>0$, tetapkan
$a=1-{1\over 4}\epsilon^2$, sehingga $\sqrt{1-a}=\bover{\epsilon}2$.
Maka terdapat $n_0\in\Bbb N$ sedemikian sehingga
$q_n(x)-\sqrt{1-x}\le\bover{\epsilon}2$ untuk setiap $x\in[0,a]$ dan
$n\ge n_0$. Khususnya, $q_n(a)\le\epsilon$, sehingga $q_n(x)\le\epsilon$
dan $q_n(x)-\sqrt{1-x}\le\epsilon$ jika $x\in[a,1]$ dan $n\ge n_0$.
Ini berarti bahwa
      
\Centerline{$0\le q_n(x)-\sqrt{1-x}\le\epsilon
\Forall n\ge n_0,\,x\in[0,1]$;}
      
\noindent karena $\epsilon$ sebarang,
$\sequencen{q_n(x)}\to\sqrt{1-x}$ secara seragam pada $[0,1]$.
      
\medskip
      
{\bf (b)} Sekarang tetapkan $p_n(x)=q_n(1-x^2)$ untuk $x\in\Bbb R$.
Karena setiap $q_n$ adalah polinom real berderajat $n$, setiap $p_n$ adalah
polinom real berderajat $2n$. Selanjutnya,
      
$$\eqalign{\sup_{|x|\le 1}|p_n(x)-|x||
&=\sup_{|x|\le 1}|q_n(1-x^2)-\sqrt{1-(1-x^2)}|\cr
&=\sup_{y\in[0,1]}|q_n(y)-\sqrt{1-y}|
\to 0\cr}$$
      
\noindent ketika $n\to\infty$, sehingga
$\lim_{n\to\infty}p_n(x)=|x|$ secara seragam untuk $|x|\le 1$,
sebagaimana diperlukan.
}%end of proof of 281C
      
\leader{281D}{Akibat} Misalkan $X$ suatu himpunan dan $A$ subruang linear
tertutup dalam norma dari $\ell^{\infty}(X)$ yang memuat $\chi X$ dan memenuhi
$f\times g\in A$ setiap kali $f$, $g\in A$. Maka $|f|\in A$ untuk setiap
$f\in A$.
      
\proof{ Tetapkan
      
\Centerline{$f_1=\Bover1{1+\|f\|_{\infty}}f$,}
      
\noindent sehingga $f_1\in A$ dan $\|f_1\|_{\infty}\le 1$. Karena $A$ memuat
$\chi X$ dan tertutup terhadap perkalian,
$p\frsmallcirc f_1\in A$ untuk setiap polinom $p$ dengan koefisien real.
Khususnya, $g_n=p_n\frsmallcirc f_1\in A$ untuk setiap $n$, dengan
$\sequencen{p_n}$ adalah barisan dari 281C. Kini, karena $|f_1(x)|\le 1$
untuk setiap $x\in X$,
      
\Centerline{$\|g_n-|f_1|\|_{\infty}
=\sup_{x\in X}|p_n(f_1(x))-|f_1(x)||
\le\sup_{|y|\le 1}|p_n(y)-|y||
\to 0$}
      
\noindent ketika $n\to\infty$. Karena $A$ tertutup dalam norma
$\|\,\|_{\infty}$, $|f_1|\in A$; akibatnya $|f|\in A$, sebagaimana diklaim.
}%end of proof of 281D
\leader{281E}{Teorema Stone-Weierstrass: bentuk kedua} Misalkan $X$ adalah
ruang topologis dan $K$ subhimpunan kompak dari $X$. Tuliskan $C_b(X)$ untuk
ruang semua fungsi kontinu bernilai real yang terbatas pada $X$. Misalkan
$A\subseteq C_b(X)$ sedemikian sehingga
      
\inset{$A$ adalah subruang linear dari $C_b(X)$;}
      
\inset{$f\times g\in A$ untuk setiap $f$, $g\in A$;}
      
\inset{$\chi X\in A$;}
      
\inset{untuk setiap titik berbeda $x$, $y$ dari $K$ terdapat $f\in A$
sedemikian sehingga $f(x)\ne f(y)$.}
      
\noindent Maka untuk setiap $h:K\to\Bbb R$ kontinu dan $\epsilon>0$
terdapat $f\in A$ sedemikian sehingga
      
\inset{$|f(x)-h(x)|\le\epsilon$ untuk setiap $x\in K$,}
      
\inset{jika $K\ne\emptyset$, $\inf_{x\in X}f(x)\ge\inf_{x\in K}h(x)$ dan
$\sup_{x\in X}f(x)\le\sup_{x\in K}h(x)$.}
      
\proof{ Misalkan $\overline{A}$ adalah penutupan dalam norma
$\|\,\|_{\infty}$ dari $A$ di $\ell^{\infty}(X)$. Bermanfaat untuk mengetahui bahwa
$\overline{A}\subseteq C_b(X)$; ini karena limit seragam fungsi-fungsi kontinu
adalah kontinu. (Namun jika hal ini baru bagi Anda, atau ingatan Anda telah
memudar, jangan berhenti untuk mencarinya sekarang; cukup baca
`$\overline{A}\cap C_b(X)$' menggantikan `$\overline{A}$' dalam sisa argumen
ini.) Menurut 281B--281D, $\overline{A}$ adalah subruang linear dari $C_b(X)$
dan $|f|\in\overline{A}$ untuk setiap $f\in\overline{A}$, sehingga kondisi
281A berlaku untuk $\overline{A}$.
      
Ambil fungsi kontinu $h:K\to\Bbb R$ dan $\epsilon>0$. Kasus $K=\emptyset$
atau $h$ konstan bersifat trivial, karena semua fungsi konstan termasuk dalam
$A$; jadi anggaplah $M_0=\inf_{x\in K}h(x)$ dan
$M_1=\sup_{x\in K}h(x)$ terdefinisi dan berbeda. Seperti diamati di akhir
bukti 281A, $M_0$ dan $M_1$ berhingga. Tetapkan
      
\Centerline{$\eta=\min(\bover13\epsilon,\bover12(M_1-M_2))>0$,\quad
$\tilde h(x)=\med(M_0+\eta,h(x),M_1-\eta)$ untuk $x\in K$}
      
\noindent (definisi: 2A1Ac), sehingga $\tilde h:K\to\Bbb R$ kontinu dan
$M_0+\eta\le\tilde h(x)\le M_1-\eta$ untuk setiap $x\in K$. Menurut 281A,
terdapat $f_0\in\overline{A}$ sedemikian sehingga
$|f_0(x)-\tilde h(x)|\le\eta$ untuk setiap $x\in K$ dan
$M_0+\eta\le f_0(x)\le M_1-\eta$ untuk setiap $x\in X$. Kini terdapat
$f\in A$ sedemikian sehingga $\|f-f_0\|_{\infty}\le\eta$, sehingga
      
\Centerline{$|f(x)-h(x)|\le|f(x)-f_0(x)|+|f_0(x)-\tilde h(x)|
+|\tilde h(x)-h(x)|\le 3\eta\le\epsilon$}
      
\noindent untuk setiap $x\in K$, sedangkan
      
\Centerline{$M_0\le f_0(x)-\eta\le f(x)\le f_0(x)+\eta\le M_1$}
      
\noindent untuk setiap $x\in X$.
}%end of proof of 281E
      
\leader{281F}{Akibat: teorema Weierstrass} Misalkan $K$ sebarang subhimpunan
tertutup dan terbatas dari $\Bbb R$. Maka setiap $h:K\to\Bbb R$ kontinu dapat
dihampiri secara seragam pada $K$ oleh polinom.
      
\proof{ Terapkan 281E dengan $X=K$ (perhatikan bahwa $K$, karena tertutup
dan terbatas, adalah kompak), dan $A$ himpunan polinom berkoefisien real,
yang dipandang sebagai fungsi dari $K$ ke $\Bbb R$.
}%end of proof of 281F
      
\leader{281G}{Teorema Stone-Weierstrass: bentuk ketiga} Misalkan $X$ adalah
ruang topologis dan $K$ subhimpunan kompak dari $X$. Tuliskan $C_b(X;\Bbb C)$
untuk ruang semua fungsi kontinu bernilai kompleks yang terbatas pada $X$,
sehingga $C_b(X;\Bbb C)$ adalah ruang linear atas $\Bbb C$. Misalkan
$A\subseteq C_b(X;\Bbb C)$ sedemikian sehingga
      
\inset{$A$ adalah subruang linear dari $C_b(X;\Bbb C)$;}
      
\inset{$f\times g\in A$ untuk setiap $f$, $g\in A$;}
      
\inset{$\chi X\in A$;}
      
\inset{konjugat kompleks $\bar f$ dari $f$ termasuk dalam $A$ untuk setiap
$f\in A$;}
      
\inset{untuk setiap titik berbeda $x$, $y$ dari $K$ terdapat $f\in A$
sedemikian sehingga $f(x)\ne f(y)$.}
      
\noindent Maka untuk setiap $h:K\to\Bbb C$ kontinu dan $\epsilon>0$
terdapat $f\in A$ sedemikian sehingga
      
\inset{$|f(x)-h(x)|\le\epsilon$ untuk setiap $x\in K$,}
      
\inset{jika $K\ne\emptyset$, $\sup_{x\in X}|f(x)|\le\sup_{x\in K}|h(x)|$.}
      
\proof{ Jika $K=\emptyset$, atau $h$ identik nol, kita dapat mengambil
$f=\tbf{0}$. Jadi anggaplah $M=\sup_{x\in K}|h(x)|>0$.
      
\medskip
      
{\bf (a)} Tetapkan
      
\Centerline{$A_{\Bbb R}=\{f:f\in A,\,f(x)$ real untuk setiap $x\in X\}$.}
      
\noindent Maka $A_{\Bbb R}$ memenuhi kondisi 281E. \Prf\ (i) Jelas
$A_{\Bbb R}$ adalah subhimpunan dari $C_b(X)=C_b(X;\Bbb R)$, tertutup
terhadap penjumlahan, perkalian dengan skalar real, dan perkalian titik demi
titik fungsi, serta memuat $\chi X$. Jika $x$, $y$ adalah titik berbeda dari
$K$, terdapat $f\in A$ sedemikian sehingga $f(x)\ne f(y)$. Kini
      
\Centerline{$\Real f=\Bover12(f+\bar f)$,
\quad$\Imag f=\Bover1{2i}(f-\bar f)$}
      
\noindent keduanya termasuk dalam $A$ dan bernilai real, sehingga termasuk
dalam $A_{\Bbb R}$; dan setidaknya salah satunya mengambil nilai berbeda di
$x$ dan $y$. \Qed
      
\medskip
      
{\bf (b)} Akibatnya, untuk fungsi kontinu $h:K\to\Bbb C$ dan $\epsilon>0$,
kita dapat menerapkan 281E dua kali untuk memperoleh $f_1$, $f_2\in A_{\Bbb R}$
sedemikian sehingga
      
\Centerline{$|f_1(x)-\Real(h(x))|\le\eta$,\quad
$|f_2(x)-\Imag(h(x))|\le\eta$}
      
\noindent untuk setiap $x\in K$, dengan
$\eta=\min(\bover12,M,\bover16\epsilon)>0$. Dengan menetapkan
$g=f_1+if_2$, kita memiliki $g\in A$ dan
$|g(x)-h(x)|\le2\eta$ untuk setiap $x\in K$.
      
\medskip
      
{\bf (c)} Tetapkan $M_1=\|g\|_{\infty}$. Jika $M_1\le M$, kita dapat
mengambil $f=g$ dan berhenti. Jika tidak, pertimbangkan fungsi
      
\Centerline{$\phi(t)=\Bover{M-\eta}{\max(M,\sqrt{t})}$}
%\eta\le M so \phi\ge 0
      
\noindent untuk $t\in [0,M_1^2]$. Menurut teorema Weierstrass (281F),
terdapat polinom real $p$ sedemikian sehingga
$|\phi(t)-p(t)|\le\Bover{\eta}{M_1}$ setiap kali $0\le t\le M_1^2$.
Perhatikan bahwa $|g|^2=g\times\bar g\in A$, sehingga
      
\Centerline{$f=g\times p(|g|^2)\in A$.}
      
\noindent Sekarang
      
\Centerline{$|p(t)|\le\phi(t)+\Bover{\eta}{M_1}
\le\phi(t)+\Bover{\eta}{\max(M,\sqrt{t})}
=\Bover{M}{\max(M,\sqrt{t})}$}
      
\noindent setiap kali $0\le t\le M_1^2$, sehingga
      
\Centerline{$|f(x)|\le|g(x)|\Bover{M}{\max(M,|g(x)|)}\le M$}
      
\noindent untuk setiap $x\in X$. Selanjutnya, jika
$0\le t\le\min(M_1,M+2\eta)^2$,
      
\Centerline{$|1-p(t)|\le\Bover{\eta}{M_1}+1-\phi(t)
\le\Bover{\eta}M+1-\Bover{M-\eta}{M+2\eta}
\le\Bover{4\eta}M$.}
      
\noindent Akibatnya, jika $x\in K$, maka
      
\Centerline{$|g(x)|\le\min(M_1,|h(x)|+2\eta)\le
\min(M_1,M+2\eta)$,}
      
\noindent sehingga
      
\Centerline{$|1-p(|g(x)|^2)|\le\Bover{4\eta}M$,}
      
\noindent dan
      
$$\eqalign{|f(x)-h(x)|
&\le|g(x)-h(x)|+|g(x)||1-p(|g(x)|^2)|\cr
&\le 2\eta+\Bover{4\eta}M(M+2\eta)
\le 2\eta+\Bover{4\eta}M(M+1)
%\eta\le\bover12
\le\epsilon,\cr}$$
%\eta\le\bover16\epsilon
      
\noindent sebagaimana diperlukan.
}% end of proof of 281G
      
\cmmnt{
\medskip
      
\noindent{\bf Catatan} Tentu saja kita dapat menghemat usaha dengan menerima
      
\Centerline{$\sup_{x\in X}|f(x)|\le 2\sup_{x\in K}|h(x)|$,}
      
\noindent yang sudah cukup baik untuk penerapan-penerapan di bawah.
}
      
\leader{281H}{Akibat} Misalkan $[a,b]\subseteq\Bbb R$ adalah interval tertutup
tak kosong yang terbatas dan $h:[a,b]\to\Bbb C$ fungsi kontinu. Maka untuk
setiap $\epsilon>0$ terdapat $y_0,\ldots,y_n\in\Bbb R$ dan
$c_0,\ldots,c_n\in\Bbb C$ sedemikian sehingga
      
\Centerline{$|h(x)-\sum_{k=0}^nc_ke^{iy_kx}|\le\epsilon$ untuk setiap
$x\in[a,b]$,}
      
\Centerline{$\sup_{x\in\Bbb
R}|\sum_{k=0}^nc_ke^{iy_kx}|\le\sup_{x\in[a,b]}|h(x)|$.}
      
\proof{ Terapkan 281G dengan $X=\Bbb R$, $K=[a,b]$, dan $A$ subruang linear
yang dibentang oleh fungsi-fungsi $x\mapsto e^{iyx}$ ketika $y$ menjelajahi
$\Bbb R$.}
      
\leader{281I}{Akibat} Misalkan $S^1$ adalah lingkaran satuan
$\{z:|z|=1\}\subseteq\Bbb C$. Maka untuk setiap fungsi kontinu
$h:S^1\to\Bbb C$ dan $\epsilon>0$, terdapat $n\in\Bbb N$ dan
$c_{-n},c_{-n+1},\ldots,c_0,\ldots,c_n\in\Bbb C$ sedemikian sehingga
$|h(z)-\sum_{k=-n}^nc_kz^k|\le\epsilon$ untuk setiap $z\in S^1$.
      
\proof{ Terapkan 281G dengan $X=K=S^1$ dan $A$ subruang linear yang dibentang
oleh fungsi-fungsi $z\mapsto z^k$ untuk $k\in\Bbb Z$.
}
      
\leader{281J}{Akibat} Misalkan $h:[-\pi,\pi]\to\Bbb C$ adalah fungsi kontinu
sedemikian sehingga $h(\pi)=h(-\pi)$. Maka untuk setiap $\epsilon>0$
terdapat $n\in\Bbb N$, $c_{-n},\ldots,c_n\in\Bbb C$ sedemikian sehingga
$|h(x)-\sum_{k=-n}^nc_ke^{ikx}|\le\epsilon$ untuk setiap
$x\in[-\pi,\pi]$.
      
\proof{ Intinya adalah bahwa $\tilde h:S^1\to\Bbb C$ kontinu pada $S^1$,
dengan $\tilde h(z)=h(\arg z)$; hal ini karena $\arg$ kontinu di mana-mana
kecuali di $-1$, dan
      
\Centerline{$\lim_{x\downarrow-\pi}h(x)
=h(-\pi)=h(\pi)=\lim_{x\uparrow\pi}h(x)$,}
      
\noindent sehingga
      
\Centerline{$\lim_{z\in S^1,z\to-1}\tilde h(z)=h(\pi)=\tilde h(-1)$.}
      
\noindent Sekarang menurut 281I terdapat $c_{-n},\ldots,c_n\in\Bbb C$
sedemikian sehingga $|\tilde h(z)-\sum_{k=-n}^nc_kz^k|\le\epsilon$ untuk
setiap $z\in S^1$, dan koefisien-koefisien ini juga berlaku untuk $h$.
}
\vleader{60pt}{281K}{Akibat} Misalkan $r\ge 1$ dan $K\subseteq\BbbR^r$
adalah himpunan tertutup terbatas tak kosong. Misalkan $h:K\to\Bbb C$ adalah
fungsi kontinu dan $\epsilon>0$. Maka terdapat $y_0,\ldots,y_n\in\Bbb Q^r$
dan $c_0,\ldots,c_n\in\Bbb C$ sedemikian sehingga
      
\Centerline{$|h(x)-\sum_{k=0}^nc_ke^{iy_k\dotproduct x}|\le\epsilon$ untuk
setiap $x\in K$,}
      
\Centerline{$\sup_{x\in\Bbb R^r}|\sum_{k=0}^nc_ke^{iy_k\dotproduct x}|
\le\sup_{x\in K}|h(x)|$,}
      
\noindent dengan menuliskan $y\dotproduct x=\sum_{j=1}^r\eta_j\xi_j$ ketika
$y=(\eta_1,\ldots,\eta_r)$ dan $x=(\xi_1,\ldots,\xi_r)$ termasuk dalam
$\BbbR^r$.
      
      
\proof{ Terapkan 281G dengan $X=\BbbR^r$ dan $A$ subruang linear yang
dibentang oleh fungsi-fungsi $x\mapsto e^{iy\dotproduct x}$ ketika $y$
menjelajahi $\BbbQ^r$.}
      
\leader{281L}{Akibat} Misalkan $r\ge 1$ dan $K\subseteq\BbbR^r$ adalah
himpunan tertutup terbatas tak kosong. Misalkan $h:K\to\Bbb R$ fungsi
kontinu dan $\epsilon>0$. Maka terdapat $y_0,\ldots,y_n\in\BbbR^r$ dan
$c_0,\ldots,c_n\in\Bbb C$ sedemikian sehingga, jika
$g(x)=\sum_{k=0}^nc_ke^{iy_k\dotproduct x}$, maka $g$ bernilai real dan
      
\Centerline{$|h(x)-g(x)|\le\epsilon$ untuk setiap $x\in K$,}
      
\Centerline{$\inf_{y\in K}h(y)\le g(x)\le\sup_{y\in K}h(y)$ untuk setiap
$x\in\BbbR^r$.}
      
\proof{ Terapkan 281E dengan $X=\BbbR^r$ dan $A$ himpunan fungsi {\it bernilai
real} pada $\BbbR^r$ yang merupakan kombinasi {\it linear kompleks} dari fungsi-
fungsi $x\mapsto e^{iy\dotproduct x}$; seperti dicatat pada bagian (a) bukti
281G, $A$ memenuhi kondisi 281E.}
      
\leader{281M}{Teorema ekuidistribusi Weyl}\cmmnt{ Kita kini siap untuk
salah satu hasil dasar teori bilangan. Sebenarnya saya akan menerapkannya
untuk memberikan contoh dalam \S285 di bawah, tetapi (setidaknya pada kasus
satu variabel) teorema ini pasti termasuk dalam daftar (yang agak panjang)
hal-hal yang patut diketahui setiap matematikawan murni. Demi penerapan yang
saya maksud, saya memberikan versi lengkap berdimensi $r$, tetapi Anda mungkin
ingin memahaminya terlebih dahulu dengan $r=1$.
      
Akan bermanfaat memiliki notasi untuk `bagian pecahan'.} Untuk setiap bilangan
real $x$, tuliskan $\fraction{x}$ untuk bilangan dalam $\coint{0,1}$ sedemikian
sehingga $x-\fraction{x}$ adalah bilangan bulat. \cmmnt{Sekarang untuk
teoremanya.}
      
\leader{281N}{Teorema} Misalkan $\eta_1,\ldots,\eta_r$ bilangan real sedemikian
sehingga $1,\eta_1,\ldots,\eta_r$ bebas linear atas $\Bbb Q$. Maka setiap kali
$0\le \alpha_j\le\beta_j\le 1$ untuk setiap $j\le r$,
      
\Centerline{$\lim_{n\to\infty}\Bover1{n+1}\#(\{m:m\le
n,\,\fraction{m\eta_j}\in[\alpha_j,\beta_j]$ untuk setiap $j\le r\})
=\prod_{j=1}^r(\beta_j-\alpha_j)$.}
      
\cmmnt{\medskip
      
\noindent{\bf Catatan} Jadi teorema ini menyatakan bahwa proporsi jangka
panjang $r$-tuple $(\fraction{m\eta_1},
\ifdim\pagewidth=390pt\break\fi
\ldots,\fraction{m\eta_r})$ yang termasuk dalam interval
$[a,b]\subseteq[\tbf{0},\tbf{1}]$ tepat sama dengan ukuran Lebesgue $\mu[a,b]$
dari interval tersebut. Tentu saja kondisi `$1,\eta_1,\ldots,\eta_r$ bebas
linear atas $\Bbb Q$' diperlukan sekaligus cukup (281Xg).}
      
\proof{{\bf (a)} Tuliskan $y=(\eta_1,\ldots,\eta_r)\in\BbbR^r$,
      
\Centerline{$\fraction{my}
=(\fraction{m\eta_1},\ldots,\fraction{m\eta_r})
\in\coint{\tbf{0},\tbf{1}}=\coint{0,1}^r$}
      
\noindent untuk setiap $m\in\Bbb N$. Tetapkan $I=[\tbf{0},\tbf{1}]=[0,1]^r$,
dan untuk setiap fungsi $f:I\to\Bbb R$ tuliskan
      
\Centerline{$\overline{L}(f)
=\limsup_{n\to\infty}\Bover1{n+1}\sum_{m=0}^{n}f(\fraction{my})$,}
      
\Centerline{$\underline{L}(f)
=\liminf_{n\to\infty}\Bover1{n+1}\sum_{m=0}^{n}f(\fraction{my})$;}
      
\noindent dan untuk $f:I\to\Bbb C$ tuliskan
      
\Centerline{$L(f)
=\lim_{n\to\infty}\Bover1{n+1}\sum_{m=0}^nf(\fraction{my})$}
      
\noindent jika limit tersebut ada. Patut dicatat bahwa untuk fungsi-fungsi
nonnegatif $f$, $g$, $h:I\to\Bbb R$ dengan $h\le f+g$,
      
\Centerline{$\overline{L}(h)\le\overline{L}(f)+\overline{L}(g)$,}
      
\noindent dan bahwa $L(cf+g)=cL(f)+L(g)$ untuk dua fungsi $f$, $g:I\to\Bbb C$
dengan $L(f)$ dan $L(g)$ keduanya ada, serta setiap $c\in\Bbb C$.
      
\medskip
      
{\bf (b)} Saya bermaksud menunjukkan bahwa $L(f)$ ada dan sama dengan
$\int_{I}f$ untuk (banyak) fungsi kontinu $f$. Langkah kuncinya adalah
mempertimbangkan fungsi-fungsi berbentuk
      
\Centerline{$f(x)=e^{2\pi ik\dotproduct x}$,}
      
\noindent dengan $k=(\kappa_1,\ldots,\kappa_r)\in\Bbb Z^r$. Dalam kasus ini,
jika $k\ne\tbf{0}$,
      
\Centerline{$k\dotproduct y=\sum_{j=1}^r\kappa_j\eta_j\notin\Bbb Z$}
      
\noindent karena $1,\eta_1,\ldots,\eta_r$ bebas linear atas $\Bbb Q$. Jadi
      
$$\eqalignno{L(f)
&=\lim_{n\to\infty}\Bover1{n+1}
  \sum_{m=0}^ne^{2\pi ik\dotproduct \fraction{my}}
=\lim_{n\to\infty}\Bover1{n+1}\sum_{m=0}^ne^{2\pi imk\dotproduct y}\cr
\noalign{\noindent (karena
$mk\dotproduct y-k\dotproduct \fraction{my}
=\sum_{j=1}^r\kappa_j(m\eta_j-\fraction{m\eta_j})$ adalah bilangan bulat)}
&=\lim_{n\to\infty}
\bover{1-e^{2\pi i(n+1)k\dotproduct y}}{(n+1)(1-e^{2\pi ik\dotproduct
y})}\cr
\noalign{\noindent (karena $e^{2\pi ik\dotproduct y}\ne 1$)}
&=0,\cr}$$
      
\noindent karena $|1-e^{2\pi i(n+1)k\dotproduct y}|\le 2$ untuk setiap $n$.
Tentu saja kita juga dapat menghitung integral $f$ pada $I$, yaitu
      
$$\eqalignno{\int_If(x)dx
&=\int_Ie^{2\pi ik\dotproduct x}dx
=\int_I\prod_{j=1}^re^{2\pi i\kappa_j\xi_j}dx\cr
\noalign{\noindent (dengan menuliskan $x=(\xi_1,\ldots,\xi_r)$)}
&=\int_{0}^1\ldots\int_0^1\prod_{j=1}^r
  e^{2\pi i\kappa_j\xi_j}d\xi_r\ldots d\xi_1\cr
&=\int_{0}^1e^{2\pi i\kappa_r\xi_r}d\xi_r
  \ldots\int_0^1e^{2\pi i\kappa_1\xi_1}d\xi_1
=0\cr}$$
      
\noindent karena setidaknya satu $\kappa_j$ tidak nol, dan untuk $j$ tersebut
kita harus memiliki
      
\Centerline{$\int_0^1e^{2\pi i\kappa_j\xi_j}d\xi_j
=\Bover{1}{2\pi i\kappa_j}(e^{2\pi i\kappa_j}-1)=0$.}
      
\noindent Jadi $L(f)=\int_If=0$ ketika $k\ne\tbf{0}$. Sebaliknya, jika
$k=\tbf{0}$, maka $f$ konstan dengan nilai $1$, sehingga
      
\Centerline{$L(f)
=\lim_{n\to\infty}\Bover1{n+1}\sum_{m=0}^nf(\fraction{my})
=\lim_{n\to\infty}1=1=\int_If(x)dx$.}
      
\medskip
      
{\bf (c)} Sekarang tuliskan
$\partial I=[\tbf{0},\tbf{1}]\setminus\ooint{\tbf{0},\tbf{1}}$, yaitu batas
dari $I$. Jika $f:I\to\Bbb C$ kontinu dan $f(x)=0$ untuk $x\in\partial I$,
maka $L(f)=\int_If$. \Prf\ Seperti dalam 281I, misalkan $S^1$ adalah
lingkaran satuan $\{z:z\in\Bbb C,\,|z|=1\}$, dan tetapkan
$K=(S^1)^r\subseteq\Bbb C^r$. Jika kita memandang $K$ sebagai subhimpunan
dari $\BbbR^{2r}$, maka himpunan itu tertutup dan terbatas. Misalkan
$\phi:K\to I$ diberikan oleh
      
\Centerline{$\phi(\zeta_1,\ldots,\zeta_r)
=(\Bover12+\Bover{\arg\zeta_1}{2\pi},\ldots,
\Bover12+\Bover{\arg\zeta_r}{2\pi})$}
      
\noindent untuk $\zeta_1,\ldots,\zeta_r\in S^1$. Maka $h=f\phi:K\to\Bbb
C$ kontinu, karena $\phi$ kontinu pada $(S^1\setminus\{-1\})^r$ dan
      
\Centerline{$\lim_{w\to z}f\phi(w)=f\phi(z)=0$}
      
\noindent untuk setiap $z\in K\setminus(S^1\setminus\{-1\})^r$.
(Bandingkan 281J.) Sekarang terapkan 281G dengan $X=K$ dan $A$ himpunan
polinom dalam $\zeta_1,\ldots,\zeta_r,\zeta_1^{-1},\ldots,\zeta_r^{-1}$
untuk melihat bahwa, diberikan $\epsilon>0$, terdapat fungsi berbentuk
      
\Centerline{$g(z)=\sum_{k\in
J}c_k\zeta_1^{\kappa_1}\ldots\zeta_r^{\kappa_r}$,}
      
\noindent untuk suatu himpunan berhingga $J\subseteq\Bbb Z^r$ dan konstanta
$c_k\in\Bbb C$ untuk $k\in J$, sedemikian sehingga
      
\Centerline{$|g(z)-h(z)|\le\epsilon$ untuk setiap $z\in K$.}
      
\noindent Tetapkan
      
\Centerline{$\tilde g(x)
=g(e^{\pi i(2\xi_1-1)},\ldots,e^{\pi i(2\xi_r-1)})
=\sum_{k\in J}c_ke^{\pi ik\dotproduct (2x-\tbf{1})}
=\sum_{k\in J}(-1)^{k\tbf{.1}}c_ke^{2\pi ik\dotproduct x}$,}
      
\noindent sehingga $\tilde g\phi=g$, dan perhatikan bahwa
      
\Centerline{$\sup_{x\in I}|\tilde g(x)-f(x)|
=\sup_{z\in K}|g(z)-h(z)|\le\epsilon$.}
      
Kini $\tilde g$ berbentuk seperti yang ditangani pada (a), sehingga
$L(\tilde g)=\int_I\tilde g$. Ambil $n_0$ sedemikian sehingga
      
\Centerline{$\bigl|\int_I\tilde g
  -\Bover1{n+1}\sum_{m=0}^n\tilde g(\fraction{my})\bigr|
\le\epsilon$}
      
\noindent untuk setiap $n\ge n_0$. Maka
      
\Centerline{$|\int_If-\int_I\tilde g|\le\int_I|f-\tilde g|\le\epsilon$}
      
\noindent dan
      
$$\eqalign{|\Bover1{n+1}\sum_{m=0}^{n}\tilde g(\fraction{my})
 -\Bover1{n+1}\sum_{m=0}^{n}f(\fraction{my})|
&\le\Bover1{n+1}\sum_{m=0}^{n}
 |\tilde g(\fraction{my})-f(\fraction{my})|\cr
&\le\Bover1{n+1}(n+1)\epsilon
=\epsilon\cr}$$
      
\noindent untuk setiap $n\in\Bbb N$. Jadi untuk $n\ge n_0$ kita harus
memiliki
      
\Centerline{$|\Bover1{n+1}\sum_{m=0}^{n}f(\fraction{my})
 -\int_If|\le3\epsilon$.}
      
\noindent Karena $\epsilon$ sebarang, $L(f)=\int_If$, sebagaimana diperlukan.
\Qed
      
\medskip
      
{\bf (d)} Perhatikan selanjutnya bahwa jika $a$,
$b\in\ooint{\tbf{0},\tbf{1}}=\ooint{0,1}^r$, dan $\epsilon>0$, terdapat
fungsi kontinu $f_1$, $f_2$ sedemikian sehingga
      
\Centerline{$f_1\le\chi[a,b]\le f_2\le\chi\ooint{\tbf{0},\tbf{1}}$,
\quad $\int_If_2-\int_If_1\le\epsilon$.}
      
\noindent\Prf\ Ini elementer. Untuk $n\in\Bbb N$, definisikan
$h_n:\Bbb R\to[0,1]$ dengan menetapkan $h_n(\xi)=0$ jika $\xi\le 0$,
$2^n\xi$ jika $0\le\xi\le 2^{-n}$, dan $1$ jika $\xi\ge 2^{-n}$.
Tetapkan
      
\Centerline{$f_{1n}(x)
=\prod_{j=1}^rh_n(\xi_j-\alpha_j)h_n(\beta_j-\xi_j)$,}
      
\Centerline{$f_{2n}(x)
=\prod_{j=1}^r(1-h_n(\alpha_j-\xi_j))(1-h_n(\xi_j-\beta_j))$}
      
\noindent untuk $x=(\xi_1,\ldots,\xi_r)\in\BbbR^r$. (Bandingkan bukti
242Oa.) Maka $f_{1n}\le\chi[a,b]\le f_{2n}$ untuk setiap $n$,
$f_{2n}\le\chi\ooint{\tbf{0},\tbf{1}}$ untuk semua $n$ yang cukup besar
sehingga
      
\Centerline{$2^{-n}
\le\min(\min_{j\le r}\alpha_j,\min_{j\le r}(1-\beta_j))$,}
      
\noindent dan $\lim_{n\to\infty}f_{2n}(x)-f_{1n}(x)=0$ untuk setiap $x$,
sehingga
      
\Centerline{$\lim_{n\to\infty}\int_If_{2n}-\int_If_{1n}=0$.}
      
\noindent Jadi kita dapat mengambil $f_1=f_{1n}$, $f_2=f_{2n}$ untuk setiap
$n$ yang cukup besar. \Qed
      
\medskip
      
{\bf (e)} Dari sini, jika $a$, $b\in\ooint{\tbf{0},\tbf{1}}$ dan $a\le b$,
$L(\chi[a,b])=\mu[a,b]$. \Prf\ Ambil $\epsilon>0$. Ambil $f_1$, $f_2$
seperti pada (d). Kemudian, menggunakan (c),
      
\Centerline{$\overline{L}(\chi[a,b])\le\overline{L}(f_2)
=L(f_2)=\int_If_2\le\int_If_1+\epsilon\le\mu[a,b]+\epsilon$,}
      
\Centerline{$\underline{L}(\chi[a,b])\ge\underline{L}(f_1)
=L(f_1)=\int_If_1\ge\int_If_2-\epsilon\ge\mu[a,b]-\epsilon$,}
      
\noindent sehingga
      
\Centerline{$\mu[a,b]-\epsilon\le\underline{L}(\chi[a,b])
\le\overline{L}(\chi[a,b])\le\mu[a,b]+\epsilon$.}
      
\noindent Karena $\epsilon$ sebarang,
      
\Centerline{$\mu[a,b]=\overline{L}(\chi[a,b])=\underline{L}(\chi[a,b])
=L(\chi[a,b])$,}
      
\noindent sebagaimana diperlukan.\ \Qed
      
\medskip
      
{\bf (f)} Untuk menyelesaikan bukti, ambil sebarang $a$, $b\in I$ dengan
$a\le b$. Untuk $0<\epsilon\le\bover12$, tetapkan
$I_{\epsilon}=[\epsilon\tbf{1},(1-\epsilon)\tbf{1}]$, sehingga
$I_{\epsilon}$ adalah interval tertutup yang termasuk dalam
$\ooint{\tbf{0},\tbf{1}}$ dan $\mu I_{\epsilon}=(1-2\epsilon)^r$. Tentu saja
$L(\chi I)=\mu I=1$, sehingga
      
\Centerline{$L(\chi(I\setminus I_{\epsilon}))
=L(\chi I)-L(\chi I_{\epsilon})
=1-\mu I_{\epsilon}$,}
      
\noindent dan
      
$$\eqalign{\mu[a,b]-1+\mu I_{\epsilon}
&\le\mu[a,b]+\mu I_{\epsilon}-\mu([a,b]\cup I_{\epsilon})
=\mu([a,b]\cap I_{\epsilon})\cr
&=L(\chi([a,b]\cap I_{\epsilon}))
\le\underline{L}(\chi([a,b]))\cr
&\le\overline{L}(\chi([a,b]))
\le\overline{L}(\chi([a,b]\cap I_{\epsilon}))
  +\overline{L}(\chi(I\setminus I_{\epsilon}))\cr
&=L(\chi([a,b]\cap I_{\epsilon}))+1-\mu I_{\epsilon}\cr
&=\mu([a,b]\cap I_{\epsilon})+1-\mu I_{\epsilon}
\le\mu[a,b]+1-\mu I_{\epsilon}.\cr}$$
      
\noindent Karena $\epsilon$ sebarang,
      
\Centerline{$\mu[a,b]=\overline{L}(\chi[a,b])=\underline{L}(\chi[a,b])
=L(\chi[a,b])$,}
      
\noindent sebagaimana dinyatakan.
}%end of proof of 281N
\exercises{
\leader{281X}{Latihan dasar (a)}
%\spheader 281Xa
Misalkan $A$ adalah himpunan fungsi kontinu terbatas
$f:\BbbR^r\times\BbbR^r\to\Bbb R$ yang dapat dinyatakan dalam bentuk
$f(x,y)=\sum_{k=0}^ng_k(x)g'_k(y)$, dengan semua $g_k$, $g'_k$ fungsi kontinu
dari $\BbbR^r$ ke $\Bbb R$. Tunjukkan bahwa untuk setiap fungsi kontinu
terbatas $h:\BbbR^r\times\BbbR^r\to\Bbb R$, setiap himpunan terbatas
$K\subseteq\BbbR^r\times\BbbR^r$, dan setiap $\epsilon>0$, terdapat $f\in A$
sedemikian sehingga
$|f(x,y)-h(x,y)|\le\epsilon$ untuk setiap $(x,y)\in K$ dan
$\sup_{x,y\in\BbbR^r}|f(x,y)|\le\sup_{x,y\in\BbbR^r}|h(x,y)|$.
%281E
      
\spheader 281Xb Misalkan $K$ himpunan tertutup terbatas dalam $\BbbR^r$,
dengan $r\ge 1$, dan $h:K\to\Bbb R$ fungsi kontinu. Tunjukkan bahwa untuk
setiap $\epsilon>0$ terdapat polinom $p$ dalam $r$ variabel sedemikian sehingga
$|h(x)-p(x)|\le\epsilon$ untuk setiap $x\in K$.
%281F
      
\sqheader 281Xc Misalkan $[a,b]$ interval tertutup tak kosong dari $\Bbb R$
dan $h:[a,b]\to\Bbb R$ fungsi kontinu. Tunjukkan bahwa untuk setiap
$\epsilon>0$ terdapat
$y_0,\ldots,y_n,a_0,\ldots,a_n,b_0,\ldots,b_n\in\Bbb R$ sedemikian sehingga
      
\Centerline{$|h(x)-\sum_{k=0}^n(a_k\cos y_kx+b_k\sin y_kx)|\le\epsilon$
untuk setiap $x\in [a,b]$,}
      
\Centerline{$\sup_{x\in\Bbb R}|\sum_{k=0}^n(a_k\cos y_kx+b_k\sin
y_kx)|\le\sup_{x\in [a,b]}|h(x)|$.}
%281H
      
\spheader 281Xd Misalkan $h$ fungsi bernilai kompleks pada $\ocint{-\pi,\pi}$
sedemikian sehingga $|h|^p$ terintegralkan, dengan $1\le p<\infty$.
Tunjukkan bahwa untuk setiap $\epsilon>0$ terdapat fungsi berbentuk
$x\mapsto f(x)=\sum_{k=-n}^nc_ke^{ikx}$, dengan
$c_{-k},\ldots,c_k\in\Bbb C$, sedemikian sehingga
$\int_{-\pi}^{\pi}|h-f|^p\le\epsilon$. (Bandingkan 244H.)
%281H
      
\sqheader 281Xe Misalkan $h:[-\pi,\pi]\to\Bbb R$ fungsi kontinu sedemikian
sehingga $h(\pi)=h(-\pi)$, dan $\epsilon>0$. Tunjukkan bahwa terdapat
$a_0,\ldots,a_n,b_1,\ldots,b_n\in\Bbb R$ sedemikian sehingga
      
\Centerline{$|h(x)-\Bover12a_0-\sum_{k=1}^n(a_k\cos kx+b_k\sin kx)|
\le\epsilon$}
      
\noindent untuk setiap $x\in[-\pi,\pi]$.
%281J
      
\spheader 281Xf Misalkan $K$ himpunan tertutup terbatas tak kosong dalam
$\BbbR^r$, dengan $r\ge 1$, dan $h:K\to\Bbb R$ fungsi kontinu. Tunjukkan
bahwa untuk setiap $\epsilon>0$ terdapat $y_0,\ldots,y_n\in\BbbR^r$,
$a_0,\ldots,a_n,b_0,\ldots,b_n\in\Bbb R$ sedemikian sehingga
      
\Centerline{$|h(x)
  -\sum_{k=0}^n(a_k\cos(y_k\dotproduct x)+b_k\sin(y_k\dotproduct x))|
\le\epsilon$ untuk setiap $x\in K$,}
      
\Centerline{$\sup_{x\in\Bbb R}
 |\sum_{k=0}^n(a_k\cos(y_k\dotproduct x)+b_k\sin(y_k\dotproduct x))|
\le\sup_{x\in K}|h(x)|$,}
      
\noindent dengan menafsirkan $y\dotproduct x$ seperti dalam 281K.
%281K
      
\spheader 281Xg Misalkan $y_1,\ldots,y_r$ bilangan real sedemikian sehingga
$1,y_1,\ldots,y_r$ tidak bebas linear atas $\Bbb Q$. Tunjukkan bahwa terdapat
interval nontrivial $[a,b]\subseteq[\tbf{0},\tbf{1}]\subseteq\BbbR^r$
sedemikian sehingga $(\fraction{my_1},\ldots,\fraction{my_r})\notin[a,b]$
untuk setiap $m\in\Bbb Z$.
%281M
      
\spheader 281Xh Misalkan $\eta_1,\ldots,\eta_r$ bilangan real sedemikian
sehingga $1,\eta_1,\ldots,\eta_r$ bebas linear atas $\Bbb Q$. Misalkan pula
$0\le\alpha_j\le\beta_j\le 1$ untuk setiap $j\le r$. Tunjukkan bahwa untuk
setiap $\epsilon>0$ terdapat $n_0\in\Bbb N$ sedemikian sehingga
      
\Centerline{$|\prod_{j=1}^r(\beta_j-\alpha_j)
 -\Bover1{n+1}\#(\{m:k\le m\le k+n,\,\fraction{m\eta_j}
 \in[\alpha_j,\beta_j]$ untuk setiap $j\le r\})|\le\epsilon$}
      
\noindent setiap kali $n\ge n_0$ dan $k\in\Bbb N$. ({\it Petunjuk\/}:
dalam bukti 281N, tetapkan
      
\Centerline{$\overline{L}(f)
 =\limsup_{n\to\infty}\sup_{k\in\Bbb N}\Bover1{n+1}
 \sum_{m=k}^{k+n}f(\fraction{my})$.)}
%281N
      
\leader{281Y}{Latihan lanjutan (a)} Tunjukkan bahwa di bawah hipotesis 281A,
terdapat $f\in\overline{A}$, yang merupakan penutupan dalam norma
$\|\,\|_{\infty}$ dari $A$ di $C_b(X)$, sedemikian sehingga
$f\restr K=h$. ({\it Petunjuk\/}: ambil
$f=\lim_{n\to\infty}f_n$ dengan
      
\Centerline{$\|f_{n+1}-f_n\|_{\infty}\le\sup_{x\in K}|f_n(x)-h(x)|
\le 2^{-n}$}
      
\noindent untuk setiap $n\in\Bbb N$.)
%281A
      
\spheader 281Yb Misalkan $X$ ruang topologis dan $K\subseteq X$ subhimpunan
kompak. Misalkan untuk setiap titik berbeda $x$, $y$ dari $K$ terdapat fungsi
kontinu $f:X\to\Bbb R$ sedemikian sehingga $f(x)\ne f(y)$. Tunjukkan bahwa
untuk setiap $r\in\Bbb N$ dan setiap $h:K\to\BbbR^r$ kontinu terdapat
$f:X\to\BbbR^r$ kontinu yang memperluas $h$.
\Hint{pertimbangkan $r=1$ terlebih dahulu.}
%281Ya, 281A
      
\spheader 281Yc Misalkan $\langle X_i\rangle_{i\in I}$ keluarga sebarang
ruang Hausdorff kompak, dan $X$ produk mereka sebagai ruang topologis. Untuk
setiap $i$, tuliskan $C(X_i)$ untuk himpunan fungsi kontinu dari $X_i$ ke
$\Bbb R$, dan $\pi_i:X\to X_i$ untuk pemetaan koordinat. Tunjukkan bahwa
subaljabar $C(X)$ yang dibangkitkan oleh
$\{f\pi_i:i\in I,\,f\in C(X_i)\}$ rapat menurut norma $\|\,\|_{\infty}$
dalam $C(X)$. ({\it Catatan\/}: Anda perlu mengetahui bahwa $X$ kompak,
 dan bahwa jika $Z$ ruang Hausdorff kompak, maka untuk setiap $z$, $w\in Z$ yang berbeda
terdapat $f\in C(Z)$ dengan $f(z)\ne f(w)$. Untuk rujukan lihat
3A3J dan 3A3Bf dalam jilid berikutnya.)
%281A
      
\spheader 281Yd Misalkan $X$ ruang topologis dan $K$ subhimpunan kompak dari
$X$. Misalkan $A$ subruang linear dari ruang $C_b(X)$ fungsi kontinu bernilai
real yang terbatas pada $X$, sedemikian sehingga $|f|\in A$ untuk setiap
$f\in A$. Misalkan $h:K\to\Bbb R$ fungsi kontinu sedemikian sehingga untuk
setiap $x$, $y\in K$ terdapat $f\in A$ dengan $f(x)=h(x)$ dan $f(y)=h(y)$.
Tunjukkan bahwa untuk setiap $\epsilon>0$ terdapat $f\in A$ sedemikian
sehingga $|f(x)-h(x)|\le\epsilon$ untuk setiap $x\in K$.
%281A
      
\spheader 281Ye Misalkan $X$ ruang topologis kompak dan tuliskan $C(X)$ untuk
himpunan fungsi kontinu dari $X$ ke $\Bbb R$. Misalkan $h\in C(X)$, dan
$A\subseteq C(X)$ sedemikian sehingga
      
\inset{$A$ adalah subruang linear dari $C(X)$;}
      
\inset{{\it atau} $|f|\in A$ untuk setiap $f\in A$ { \it atau}
$f\times g\in A$ untuk setiap $f$, $g\in A$
{ \it atau} $f\times f\in A$ untuk setiap $f\in A$;}
      
\inset{untuk setiap $x$, $y\in X$ dan $\delta>0$ terdapat $f\in A$ sedemikian
sehingga $|f(x)-h(x)|\le\delta$ dan $|f(y)-h(y)|\le\delta$.}
      
\noindent Tunjukkan bahwa untuk setiap $\epsilon>0$ terdapat $f\in A$
sedemikian sehingga $|h(x)-f(x)|\le\epsilon$ untuk setiap $x\in X$.
%281E
      
\spheader 281Yf Misalkan $X$ ruang topologis kompak dan $A$ subruang linear
tertutup dalam norma $\|\,\|_{\infty}$ dari ruang $C(X)$ fungsi kontinu dari
$X$ ke $\Bbb R$. Tunjukkan bahwa berikut ini ekuivalen:
      
\quad (i) $|f|\in A$ untuk setiap $f\in A$;
      
\quad (ii) $f\times f\in A$ untuk setiap $f\in A$;
      
\quad (iii) $f\times g\in A$ untuk semua $f$, $g\in A$,
      
\noindent dan bahwa dalam kasus ini $A$ tertutup dalam $C(X)$ untuk topologi
yang ditentukan oleh pseudometrik
      
\Centerline{$(f,g)\mapsto|f(x)-g(x)|:
  C(X)\times C(X)\to\coint{0,\infty}$}
      
\noindent ketika $x$ menjelajahi $X$ (yaitu `topologi konvergensi titik demi
titik' pada $C(X)$).
%281E
      
\spheader 281Yg Tunjukkan bahwa di bawah hipotesis 281G terdapat
$f\in\overline{A}$, yang merupakan penutupan dalam norma $\|\,\|_{\infty}$
dari $A$ di $C_b(X;\Bbb C)$, sedemikian sehingga
$f\restr K=h$ dan (jika
$K\ne\emptyset$) $\|f\|_{\infty}=\sup_{x\in K}|h(x)|$.
%281G
      
\spheader 281Yh Misalkan $y\in\Bbb R$ irasional. Tunjukkan bahwa untuk setiap
fungsi Riemann-terintegralkan $f:[0,1]\to\Bbb R$,
      
\Centerline{$\int_0^1f(x)dx
=\lim_{n\to\infty}\Bover1{n+1}\sum_{m=0}^nf(\fraction{my})$,}
      
\noindent dengan menuliskan $\fraction{my}$ untuk bagian pecahan dari $my$.
({\it Petunjuk\/}: ingat {\it kriteria Riemann}: untuk setiap $\epsilon>0$,
terdapat $a_0,\ldots,a_n$ dengan $0=a_0\le a_1\le\ldots\le a_n=1$ dan
      
\Centerline{$\sum\{a_{j}-a_{j-1}:j\le n,\,\sup_{x\in[a_{j-1},a_j]}f(x)
 -\inf_{x\in[a_{j-1},a_j]}f(x)\ge\epsilon\}\le\epsilon$.)}
%281N
      
\ifdim\pagewidth>467pt\fontdimen3\tenrm=1.84pt
  \fontdimen4\tenrm=1.22pt\fi
\spheader 281Yi Misalkan $\sequencen{t_n}$ suatu barisan dalam $[0,1]$.
Tunjukkan bahwa berikut ini ekuivalen:
(i) $\lim_{n\to\infty}\bover1{n+1}\sum_{k=0}^nf(t_k)
\ifdim\pagewidth>467pt\penalty-100\fi
=\int_0^1f$ untuk setiap fungsi kontinu $f:[0,1]\to\Bbb R$;
(ii)
\ifdim\pagewidth=390pt\break\fi
$\lim_{n\to\infty}\bover1{n+1}\sum_{k=0}^nf(t_k)=\int_0^1f$ untuk
setiap fungsi Riemann-terintegralkan $f:[0,1]\to\Bbb R$;
(iii) $\liminf_{n\to\infty}\bover1{n+1}\#(\{k:k\le n,\,t_k\in G\})
\ge\mu G$ untuk setiap himpunan terbuka $G\subseteq[0,1]$, dengan $\mu$
ukuran Lebesgue pada $\Bbb R$;
(iv) $\lim_{n\to\infty}\bover1{n+1}\#(\{k:k\le n,\,t_k\le\alpha\})
=\alpha$ untuk setiap $\alpha\in[0,1]$;
(v) $\lim_{n\to\infty}\bover1{n+1}\#(\{k:k\le n,\,t_k\in E\})=\mu E$
untuk setiap $E\subseteq[0,1]$ sedemikian sehingga
$\mu(\interior E)=\mu\overline{E}$;
(vi) $\lim_{n\to\infty}\bover1{n+1}\sum_{k=0}^ne^{2\pi imt_k}=0$ untuk
setiap $m\ge 1$.
(Bandingkan 273J. Barisan $\sequencen{t_n}$ semacam ini disebut
{\bf terdistribusi seragam} atau {\bf tersebar merata}.)
%281N
\fontdimen3\tenrm=1.67pt\fontdimen4\tenrm=1.11pt
      
\spheader 281Yj Tunjukkan bahwa barisan $\sequencen{\fraction{\ln(n+1)}}$
tidak terdistribusi seragam.
%281Yi, 281N
      
\spheader 281Yk Berikan $[0,1]^{\Bbb N}$ ukuran produk $\lambda$. Tunjukkan
bahwa, menurut $\lambda$, hampir setiap barisan
$\sequencen{t_n}\in[0,1]^{\Bbb N}$ terdistribusi seragam dalam arti 281Yi.
\Hint{273J.}
%281Yi, 281N
      
\spheader 281Yl Misalkan $f:[0,1]^2\to\Bbb C$ fungsi kontinu. Tunjukkan bahwa
jika $\gamma\in\Bbb R$ irasional maka
$\lim_{a\to\infty}\bover1a\int_0^af(\fraction{t},\fraction{\gamma t})dt
\penalty-100=\int_{[0,1]^2}f$.
\Hint{mula-mula pertimbangkan fungsi berbentuk
$x\mapsto e^{2\pi ik\dotproduct x}$.}
%281N
 
\spheader 281Ym Suatu barisan $\sequencen{t_n}$ dalam $[0,1]$ disebut
{\bf terdistribusi dengan baik} (terhadap ukuran Lebesgue $\mu$) jika
      
\Centerline{$\liminf_{n\to\infty}\inf_{l\in\Bbb N}
  \Bover1{n+1}\#(\{k:l\le k\le l+n$, 
  $t_k\in G\})\ge\mu G$}
      
\noindent untuk setiap himpunan terbuka $G\subseteq[0,1]$.
(i) Tunjukkan bahwa $\sequencen{t_n}$ terdistribusi dengan baik jika dan hanya jika
$\lim_{n\to\infty}\sup_{l\in\Bbb N}
|\int_0^1f-\bover1{n+1}\sum_{k=l}^{l+n}f(t_k)|=0$ untuk setiap fungsi kontinu
$f:[0,1]\to\Bbb R$.
(ii) Tunjukkan bahwa $\sequencen{\fraction{n\alpha}}$ terdistribusi dengan baik
untuk setiap $\alpha$ irasional.
%281Xh 281Yi 281N
}%end of exercises
      
\endnotes{\Notesheader{281} Saya telah memberikan tiga pernyataan (281A, 281E,
dan 281G) dari teorema Stone-Weierstrass, dengan pengakuan (281F) terhadap
versi Weierstrass sendiri, dan tiga bentuk lain (281Ya, 281Yd, 281Yg) dalam
latihan. Bentuk lain lagi akan muncul dalam \S4A6 di Jilid 4. Menghadapi
keragaman ini, Anda mungkin ingin mencoba sendiri menuliskan teorema yang
mencakup sebagian atau semua versi tersebut. Saya sendiri tidak melihat cara
melakukannya tanpa menyusun daftar hipotesis dan kesimpulan alternatif yang
membingungkan. Pada titik itu saya bertanya, `sebenarnya apa itu teorema?',
dan menjawab: teorema adalah tempat berhenti dalam perjalanan kita; tempat
untuk beristirahat dan mengucapkan selamat kepada diri sendiri atas pencapaian
kita; tempat yang dapat kita pelajari untuk kenali dan gunakan sebagai titik
berangkat bagi petualangan baru; tempat yang dapat kita uraikan dan bagikan
dengan orang lain. Untuk beberapa teorema, seperti teorema terakhir Fermat,
ada pernyataan kanonik, suatu titik yang dapat ditemukan tepat. Untuk yang
lain, seperti teorema Stone-Weierstrass di sini, kita sampai pada kumpulan
hasil yang sangat berkaitan, semuanya bergantung pada susunan tertentu dari
argumen yang diuraikan dalam 281A--281G dan 281Ya (yang memperkenalkan
gagasan baru), dan semuanya berguna dengan cara yang berbeda. Memang saya
berpendapat bahwa kebanyakan penulis akan lebih menyukai versi 281Ya dan
281Yg, yang menghilangkan variabel $\epsilon$ yang muncul dalam 281A, 281E,
dan 281G, dengan mengorbankan pengambilan subruang tertutup $A$. Namun saya
menemukan bahwa akibat-akibat yang akan berguna kemudian (281H--281L) lebih
alami dinyatakan dalam istilah subruang linear yang tidak tertutup.
      
Penerapan teorema, atau teorema-teorema, atau metode ini -- pilih sendiri
ungkapan Anda -- sangat banyak; hanya beberapa yang dicantumkan di sini.
Salah satu yang tampaknya tidak berbahaya terdapat dalam 281Xa dan, dalam
varian berbeda, dalam 281Yc; keduanya sangat penting dalam ranah masing-masing.
Dalam jilid ini penerapan utama akan berupa 285L di bawah, yang bergantung pada
281K, dan mungkin patut dicatat bahwa ada pendekatan alternatif untuk hasil
khusus ini, berdasarkan gagasan dalam 282G. Namun saya menawarkan teorema
ekuidistribusi Weyl (281M--281N) sebagai bukti bahwa kita dapat berharap
menemukan kegunaan gagasan-gagasan ini di hampir setiap cabang matematika.
}
      
\discrpage
      
